\paragraph{Degrees of character quotients in characteristic $p$.}
The degree calculation on a special-fibre curve has a short proof that
also explains the extra hypothesis needed at ramified coefficient primes.

\begin{proposition}\label{geo-frobenius-character-degree}
Let $k$ be a finite field of characteristic $p$, let $C/k$ be a smooth
projective geometrically connected curve, and let $F:C\to C$ be absolute
Frobenius. Let $A$ be a finite commutative $\mathbf F_p$-algebra and put
$B=k\otimes_{\mathbf F_p}A$. Suppose vector bundles
$\mathcal H,\mathcal U,\mathcal V$ on $C$ have $A$-actions and
$A$-linear exact sequences
\begin{equation}\label{geo-hodge-conjugate-sequences}
 0\longrightarrow\mathcal U\longrightarrow\mathcal H
       \longrightarrow\mathcal V\longrightarrow0,
 \qquad
 0\longrightarrow F^*\mathcal V\longrightarrow\mathcal H
       \longrightarrow F^*\mathcal U\longrightarrow0.
\end{equation}
The actions on the Frobenius pullbacks are the pulled-back $A$-actions.
Then $\deg(e\mathcal H)=0$ for every idempotent $e\in B$.

Let $\chi:B\to k$ be a $k$-algebra character, and let $e$ be the
idempotent of its Artin local factor $B_e=eB$. If $e\mathcal H$ is
locally free as a module over $\mathcal O_C\otimes_k B_e$, then
\[
 \mathcal Q_\chi=\mathcal H\otimes_{B,\chi}k
 \quad\hbox{is a vector bundle and}\quad
                         \deg\mathcal Q_\chi=0.
\]
Here local freeness over the coefficient algebra is stronger than
local freeness merely over $\mathcal O_C$. If $B_e=k$, the stronger
condition is automatic.
\end{proposition}
\begin{proof}
Write $\sigma(c\otimes a)=c^p\otimes a$ on $B$; this acts on the
finite field coefficients and leaves $A$ fixed. It is not the
absolute Frobenius of the possibly nonreduced algebra $B$.
For a bundle $\mathcal W$ with $A$-action, pulling back the action
of $e=\sum_i c_i\otimes a_i$ raises the scalars $c_i$ to their
$p$-th powers. Hence
\[
 F^*(e\mathcal W)=\sigma(e)F^*\mathcal W.
\]
Idempotent images are direct summand bundles, so applying an
idempotent preserves both exact sequences. Frobenius is finite flat
on the smooth curve, and $\deg F^*\mathcal W=p\deg\mathcal W$.
Taking the $\sigma(e)$ part of the second sequence and the $e$ part
of the first gives
\[
 \deg\bigl(\sigma(e)\mathcal H\bigr)
       =p\{\deg(e\mathcal V)+\deg(e\mathcal U)\}
       =p\deg(e\mathcal H).
\]
If $k=\mathbf F_{p^f}$, then $\sigma^f=1$. Iteration yields
$\deg(e\mathcal H)=p^f\deg(e\mathcal H)$, proving the first claim.

The residue field of $B_e$ is $k$. A composition series of its
ideals has
$n=\operatorname{length}_{B_e}B_e=\dim_kB_e>0$ successive quotients
isomorphic to $k$. Local freeness over
$\mathcal O_C\otimes_k B_e$ makes tensoring this series with
$e\mathcal H$ exact. It gives a filtration by subbundles whose
$n$ successive quotients are each isomorphic to $\mathcal Q_\chi$.
In particular $\mathcal Q_\chi$ is locally free, and additivity of
degree gives
\[
 0=\deg(e\mathcal H)=n\deg\mathcal Q_\chi.
\]
This proves the assertion, including rank zero.
\end{proof}

For $A=\mathcal O_E/p\mathcal O_E$ and $p$ unramified in $E$,
the algebra becomes a product of copies of $k$ after a finite extension
of the constant field. The character conclusion then follows from
the two exact sequences alone. At a ramified prime the nilpotents
must be retained, and local freeness over the Artin factor supplies
their multiplicity in the proof.

The following example shows why a zero degree for the whole bundle
cannot replace that hypothesis. On $C=\mathbf P^1_k$ take
$B=k[\epsilon]/(\epsilon^2)$ and
$\mathcal H=\mathcal O(1)\oplus\mathcal O(-1)$, with
$\epsilon(u,v)=(sv,0)$ for a nonzero
$s\in H^0(C,\mathcal O(2))$. Although $\deg\mathcal H=0$, the
torsion-free quotient by the image of $\epsilon$ is
$\mathcal O(-1)$, of degree $-1$. The unsaturated quotient also
has a torsion summand of length $2$. The $B$-module is not locally
free at the zero divisor of $s$. This example concerns the inference
from total degree alone; it is not claimed to satisfy
\eqref{geo-hodge-conjugate-sequences} or to arise from the prescribed
abelian family.

To apply the proposition to the geometric construction, the two
filtrations and the required coefficient freeness must hold for its
actual integral Hodge bundle along each special-fibre curve. Its
generic character determinant must also be identified with the ratio
of the prescribed conjugate tangent determinants. These model
comparisons, including their behaviour at the boundary, remain to be
established; the abstract degree calculation does not supply them.
